\section{Discussion and Limitations}
\label{sec:discussion}

\subsection{What the numbers say}

GAME-Mal does not set a new state-of-the-art in raw accuracy on
this corpus — a well-tuned Random Forest over pruned associative
rules reaches almost identical accuracy. The contribution is
instead twofold: (i) the gate mechanism closes the gap between a
sequence model and a strong tree-based baseline on a task where
trees usually win, and (ii) it does so with a built-in,
forward-pass-native explanation mechanism that the Random Forest
cannot match without falling back to post-hoc methods. For
operational triage, a 0.5--1\% accuracy cost in exchange for
per-sample API attributions computed at inference time is a
favourable trade.

\subsection{Limitations}

\paragraph{Dataset scope.} Our evaluation is confined to eight
families of the UMD corpus. Older families (Drebin) and newer
ones (recent banking trojans) may exhibit API distributions
sufficiently different to shift the picture. A multi-corpus study
is a natural next step.

\paragraph{Sequence truncation.} Droidmon traces can exceed
$10^4$ API calls; we cap at 512 via tail-truncation. While we
verified empirically that this preserves the payload-execution
phase for the families studied, longer context would benefit
multi-stage malware and is a natural fit for recent long-context
attention variants.

\paragraph{Reproducibility of gate explanations.} Gate activations
are deterministic given a fixed model, but different random seeds
produce quantitatively different fingerprints (though the top
APIs per family are stable in our runs). Formal stability analysis
is left to future work.

\paragraph{Evasion.} A motivated attacker aware of GAME-Mal could
attempt to pad traces with benign API calls specifically crafted
to attract gate mass, displacing the true family signals. We have
not yet evaluated the model under adversarial evasion.

\paragraph{Single sandbox.} All traces were collected via a single
Droidmon-instrumented sandbox configuration. Cross-sandbox
generalisation (e.g.\ to CopperDroid or DroidScope traces) is
untested.

\subsection{Future work}

Three directions are immediately attractive. First, extending
the gate mechanism to produce per-head explanations rather than
averaging across heads — different heads appear to specialise
on different behavioural axes in preliminary inspection.
Second, combining the rule-mining and gating approaches: using
Markov-derived rules as an auxiliary input channel to the
transformer, rather than as a separate baseline. Third,
formalising the gate activations as a calibrated importance score
by matching them against ground-truth behavioural annotations for
a subset of samples.
